% !TEX root = ../algebra_02.tex
\section{Самосопряжённые операторы в евклидовом пространстве}

\begin{Definition}{Самосопряжённый оператор}{}
	Оператор $A \in \mathrm{End}(V)$ называется самосопряжённым, если $A^* = A$.
	В ортонормированном базисе матрица такого оператора симметрична ($A = A^T$).
\end{Definition}

\begin{Remark}{Связь с билинейными формами}{}
	Если оператор $A$ является самосопряжённым, то отображение $B: V \times V \to \mathbb{R}$, заданное правилом $(v, w) \mapsto (Av, w)$, является симметричной билинейной формой.
	Действительно:
	\[
	(Aw, v) = (v, Aw) = (v, A^* w) = (Av, w)
	\] 
\end{Remark}

\begin{Proposition}{Корни характеристического многочлена}{}
	Пусть $A$ — самосопряжённый оператор.
	Тогда его характеристический многочлен $\chi_A$ раскладывается на линейные множители над $\mathbb{R}$ (то есть не имеет комплексных корней с ненулевой мнимой частью).
\end{Proposition}

\begin{proof}
	Пусть $\lambda \in \mathbb{C}$ — корень характеристического многочлена $\chi_A$. Зафиксируем ОНБ $E$ и рассмотрим матрицу оператора $A = [A]_E$.
	Рассмотрим оператор умножения на матрицу $A$ в пространстве $\mathbb{C}^n$: $\mathbb{C}^n \xrightarrow{A \cdot} \mathbb{C}^n$. Его характеристический многочлен совпадает с $\chi_A$.
	Так как $\chi_A(\lambda) = 0$, существует собственный вектор $x \in \mathbb{C}^n$, $x \neq 0$, такой что $Ax = \lambda x$.
	Рассмотрим следующее произведение и транспонируем его:
	\[
	(Ax)^T \bar{x} = (\lambda x)^T \bar{x} = \lambda x^T \bar{x}
	\] 
	С другой стороны, так как матрица $A$ вещественная и симметричная ($A^T = A = \bar{A}$):
	\[
	x^T A^T \bar{x} = x^T A \bar{x} = x^T \overline{A x} = x^T \overline{\lambda x} = \bar{\lambda} x^T \bar{x}
	\] 
	Приравнивая два выражения, получаем $\lambda x^T \bar{x} = \bar{\lambda} x^T \bar{x}$.
	Пусть $x = (x_1, \dots, x_n)^T$. Так как $x \neq 0$, имеем:
	\[
	x^T \bar{x} = |x_1|^2 + \dots + |x_n|^2 \neq 0
	\] 
	Следовательно, можно сократить на $x^T \bar{x}$, откуда $\lambda = \bar{\lambda} \Rightarrow \lambda \in \mathbb{R}$.
\end{proof}

\begin{Proposition}{Ортогональность собственных векторов}{}
	Собственные векторы самосопряжённого оператора, отвечающие различным собственным значениям, ортогональны.
\end{Proposition}

\begin{proof}
	Пусть $A$ — самосопряжённый оператор, $\lambda, \mu$ — различные собственные значения ($\lambda \neq \mu$), а $v, w$ — соответствующие им собственные векторы.
	Вычислим скалярное произведение $(Av, w)$.
	С одной стороны:
	\[
	(Av, w) = (\lambda v, w) = \lambda (v, w)
	\] 
	С другой стороны:
	\[
	(Av, w) = (v, A^* w) = (v, Aw) = (v, \mu w) = \mu (v, w)
	\] 
	Приравниваем результаты:
	\[
	\lambda (v, w) = \mu (v, w) \Rightarrow (\lambda - \mu)(v, w) = 0
	\] 
	Поскольку собственные значения различны ($\lambda \neq \mu$), необходимо $(v, w) = 0$, то есть $v \perp w$.
\end{proof}

\begin{Theorem}{Спектральная теорема для самосопряжённых операторов}{}
	Пусть $V$ — евклидово пространство, $A \in \mathrm{End}(V)$ — самосопряжённый оператор.
	Тогда существует ортонормированный базис $E$, в котором матрица оператора $[A]_E$ диагональна.
\end{Theorem}

\begin{proof}
	Доказательство проведем индукцией по размерности пространства $n = \dim V$.

	\textbf{База индукции} ($n = 1$): Очевидна.
	
	\textbf{Шаг индукции} ($n - 1 \to n$):
	По доказанному ранее, характеристический многочлен имеет вещественный корень $\lambda$.
	Следовательно, у оператора $A$ существует собственный вектор $v$, отвечающий этому значению. Нормируем его: $e_1 = \frac{v}{\|v\|}$.
	Рассмотрим ортогональное дополнение $W = \mathrm{Lin}(e_1)^\perp$. Очевидно, что $\dim W = n - 1$.
	Проверим инвариантность подпространства $W$.
	Пусть $w \in W$, тогда:
	\[
	(Aw, e_1) = (w, A^* e_1) = (w, A e_1) = (w, \lambda e_1) = \lambda (w, e_1) = 0
	\] 
	Таким образом, $Aw \perp e_1$, следовательно $Aw \in W$.
	Рассмотрим сужение оператора $A_1 = A|_W \in \mathrm{End}(W)$. Очевидно, что оно также является самосопряжённым ($A_1^* = A_1$).
	По индуктивному предположению, в $(n-1)$-мерном пространстве $W$ существует ОНБ $e_2, \dots, e_n$, в котором матрица оператора $A_1$ диагональна.
	Тогда система векторов $e_1, e_2, \dots, e_n$ образует искомый ортонормированный базис всего пространства $V$.
\end{proof}
